% Part: intuitionistic-logic
% Chapter: tableaux
% Section: introduction

\documentclass[../../../include/open-logic-section]{subfiles}

\begin{document}

\olfileid[tr]{int}{tab}{int}

\olsection{Giriş}

!!^{tableau}s, belirli (aşağıya doğru dallanan) !!{signed formula}s
ağaçlarıdır; başka bir deyişle, bir doğruluk değeri işareti ($\True$ veya
$\False$) ve bir !!a{sentence} içeren
\[
\sFmla{\True}{!A} \text{ veya } \sFmla{\False}{!A}
\]
çiftlerinin ağaçlarıdır. !!^a{tableau}, birtakım \emph{varsayımlarla}
başlar. Sonraki her !!{signed formula}, çıkarım kurallarından biri uygulanarak
üretilir. Bazı çıkarım kuralları ağacın bir ucuna bir ya da daha fazla
!!{signed formula}s ekler; bazılarıysa iki yeni uç ekleyerek iki dal oluşturur.
Kuralların ürettiği !!{signed formula}s içindeki !!{formula}, kuralın
uygulandığı !!{signed formula} içindekinden daha az karmaşıktır. Bir dal hem
$\sFmla{\True}{!A}$ hem de $\sFmla{\False}{!A}$ içerdiğinde o dala
\emph{kapalı} deriz. !!a{tableau}{}daki her dal kapalıysa !!{tableau}{}nun tamamı
kapalıdır. Kapalı bir !!{tableau}, !!a{derivation} oluşturur ve bu türetim,
ilgili !!{signed formula}s kümesinin---bunlar !!{tableau}{}yu başlatmak için
kullanılmıştır---sağlanamaz olduğunu gösterir. Bu, bir $\Proves$ bağıntısı
tanımlamak için kullanılabilir:
$\Gamma \Proves !A$ ancak ve ancak $\Gamma_0 = \{!B_1,
\dots, !B_n\} \subseteq \Gamma$ olacak biçimde, şu varsayımlar için kapalı
!!{tableau} bulunan sonlu bir küme varsa:
\[
\{\sFmla{\False}{!A}, \sFmla{\True}{!B_1}, \dots, \sFmla{\True}{!B_n}\}.
\]

Sezgisel mantık için hem !!{signed formula} kavramını genişletmeli hem de
bağlaç kurallarını uyarlamalıyız. Ayrıca !!{formula}s, sezgisel !!{tableau}s
içinde bir işaretin ($\True$ veya $\False$) yanı sıra $\sigma$
\emph{öneklerine} de sahiptir. Önekler pozitif tam sayıların boş olmayan
dizileridir; yani
$\sigma \in (\PosInt)^* \setminus \{\emptyseq\}$. Bu önekleri çevrelerindeki
$\tuple{\ }$ olmadan yazar ve tek tek !!{element}s{}ı $,$ yerine $.$ ile
ayırırız. $\sigma$ bir önekse $\sigma.n$, $\sigma \concat \tuple{n}$
demektir; örneğin $\sigma = 1.2.1$ ise $\sigma.3$, $1.2.1.3$'tür.
Dolayısıyla örneğin
\[
\sFmla{\True}{!A \lif (!B \lif !C)}[1.2]
\]
bir \emph{önekli !!{signed formula}}dır (kısaca \emph{önekli
  !!{formula}}).

Sezgisel olarak önek, bir !!{formula}s{}in !!a{tableau} dalında
sağlanabileceği modeldeki dünyayı adlandırır; $\sigma$ bir dünyayı
adlandırıyorsa $\sigma.n$, $\sigma$'nın adlandırdığı dünyadan erişilebilen
bir dünyayı adlandırır.

Sezgisel modellerde erişilebilirlik bağıntısı yansımalı ve geçişlidir. Önekler
bakımından bu, $\sigma$'dan hem $\sigma$'nın kendisine hem de $\sigma$'yı
genişleten her öneke, yani $\sigma.n_1.\cdots.n_k$ biçimindeki her öneke
erişilebildiği anlamına gelir. $\sigma$'nın kendisini ve bütün genişletmelerini belirtmek
için $\sigma.*$ gösterimini kullanalım. Başka bir deyişle, $\sigma.*$
önekleri tam olarak $\sigma$'dan erişilebilen öneklerdir.

\end{document}
